\documentclass[../../../include/open-logic-section]{subfiles}

\begin{document}

\olfileid[pt]{sth}{cardinals}{cardsasords}
\olsection{Cardinais como Ordinais}

De fato, a nossa teoria dos cardinais fará uso (descarado) da nossa
teoria dos ordinais. Ou seja: definiremos os cardinais como certos
ordinais específicos. Em particular, proporemos o seguinte:

\begin{defn}\ollabel{defcardinalasordinal}
Se $A$ pode ser bem ordenado, então $\card{A}$ é o menor ordinal
$\gamma$ tal que $\cardeq{A}{\gamma}$. Para qualquer ordinal $\gamma$,
diz-se que $\gamma$ é um \emph{cardinal} sse $\gamma = \card{\gamma}$.
\end{defn}

Acabamos de usar a locução ``$A$ pode ser bem ordenado''. Como quase
sempre em matemática, a locução modal aqui é apenas um atalho para uma asserção existencial: dizer ``$A$ pode ser
bem ordenado'' é dizer ``existe uma relação que
bem ordena~$A$''.

Há contudo um problema com \olref{defcardinalasordinal}. Gostaríamos
que \emph{todo} o conjunto tivesse um tamanho, O.s.\@, que
$\card{A}$ existisse para todo~$A$. A definição que acabamos de dar,
porém, começa por uma condicional: ``\emph{Se} $A$ pode ser
bem ordenado\ldots''. Se existir algum conjunto $A$ que não possa ser
bem ordenado, a nossa definição simplesmente deixa de definir um objeto~$\card{A}$.

Logo, para usar \olref{defcardinalasordinal}, precisamos de uma garantia de que
todo o conjunto pode ser bem ordenado. Infelizmente, essa garantia não está
disponível em~$\ZF$. Assim, se queremos usar
\olref{defcardinalasordinal}, não há alternativa senão acrescentar um novo
axioma, tal como:
\begin{axiom}[Bom ordenamento]
Todo o conjunto pode ser bem ordenado.
\end{axiom}
Discutiremos se o axioma do bom ordenamento é aceitável em
\olref[choice][]{chap}. Daqui em diante, porém, simplesmente o admitiremos. E, com ele, é bastante direto provar
que os cardinais (tal como definidos na \olref{defcardinalasordinal}) existem e
comportam-se bem:

\begin{lem}\ollabel{lem:CardinalsExist}
Para todo o conjunto $A$:
\begin{enumerate}
	\item\ollabel{cardaexists} $\card{A}$ existe e é único;
	\item\ollabel{cardaapprox}  $\cardeq{\card{A}}{A}$;
	\item\ollabel{cardaidem}  $\card{A}$ é um cardinal, O.s.\@,
	$\card{A} = \card{\card{A}}$;
\end{enumerate}
\end{lem}

\begin{proof}
Fixe $A$. Pelo Bom ordenamento, existe um bom ordenamento $\tuple{A, R}$. Por
\olref[ordinals][ordtype]{thmOrdinalRepresentation}, $\tuple{A,
R}$ é isomorfo a um único ordinal, $\beta$. Logo
$\cardeq{A}{\beta}$. Por indução transfinita, existe um único
ordinal mínimo, $\gamma$, tal que $\cardeq{A}{\gamma}$. Logo $\card{A}
= \gamma$, o que estabelece \olref{cardaexists} e \olref{cardaapprox}.
Para \olref{cardaidem}, note-se que se $\delta \in \gamma$ então
$\cardless{\delta}{A}$, pela escolha de $\gamma$, logo também
$\cardless{\delta}{\gamma}$ pois a equipotência é uma relação de equivalência
(\olref[sfr][siz][equ]{equinumerosityisequi}). Logo $\gamma =
\card{\gamma}$.
%So, for reductio, suppose that there is some ordinal $\delta \in \gamma$ such that $\cardeq{\gamma}{\delta}$. Then, $\cardeq{\cardeq{A}{\gamma}}{\delta}$ so that $\cardeq{A}{\delta}$ by \olref[sfr][set][equ]{equinumerosityisequi}, which contradicts the choice of $\gamma$ as the least ordinal such that $\cardeq{A}{\gamma}$.
\end{proof}

O resultado seguinte garante o Princípio de Cantor, e mais ainda.
(Note-se que os cardinais herdam a ordem dos ordinais, O.s.\@,
$\cardfont{a} < \cardfont{b}$ sse $\cardfont{a} \in \cardfont{b}$. Ao
formular isto, usaremos letras Fraktur para objetos que sabemos serem
cardinais. Isto é bastante usual. Uma alternativa comum é usar
letras gregas, já que os cardinais são ordinais, escolhendo-as do
meio do alfabeto, por exemplo $\kappa, \lambda$.):
\begin{lem}\ollabel{lem:CardinalsBehaveRight}
Para quaisquer conjuntos $A$ e $B$:
\begin{align*}
	\cardeq{A}{B} &\text{ sse } \card{A} = \card{B}\\
	\cardle{A}{B} &\text{ sse } \card{A} \leq \card{B}\\
	\cardless{A}{B}&\text{ sse } \card{A} < \card{B}
\end{align*}
\end{lem}

\begin{proof}
Provaremos o sentido da esquerda para a direita da segunda afirmação (os
outros casos são semelhantes e ficam como exercício). Considere-se
o diagrama seguinte:
\begin{center}
	\begin{tikzpicture}
	\node (nodea) {$A$};
	\node[right = 6em of nodea] (nodeb) {$B$};
	\node[below = 2em of nodea] (nodecarda) {$\card{A}$};
	\node[below = 2em of nodeb] (nodecardb) {$\card{B}$};
	\draw[->] (nodea)--(nodeb);
	\draw[<->] (nodea)--(nodecarda);
	\draw[<->] (nodeb)--(nodecardb);
	\draw[->, dashed] (nodecarda)--(nodecardb);
\end{tikzpicture}
\end{center}
As setas duplas indicam !!{bijection}s, cuja existência é
garantida por \olref{lem:CardinalsExist}. Ao supor
$\cardle{A}{B}$, existe !!a{injection}  $A\to B$. Seguindo
as setas de $\card{A}$ para $A$, para $B$ e para $\card{B}$,
obtém-se !!a{injection} $\card{A} \to \card{B}$ (a seta tracejada).
\end{proof}\noindent Pode ainda usar-se \olref{lem:CardinalsBehaveRight}
para voltar a provar Cantor--Bernstein. Trata-se da afirmação de que se
$\cardle{A}{B}$ e $\cardle{B}{A}$ então $\cardeq{A}{B}$. A enunciamos
como \olref[sfr][siz][sb]{thm:schroder-bernstein}, mas a provamos primeiro---com algum esforço---em \olref[sfr][infinite][card-sb]{sec}.
Considere-se:

\begin{proof}[Nova prova de Schr\"oder-Bernstein]
Se $\cardle{A}{B}$ e $\cardle{B}{A}$, então $\card{A} \leq \card{B}$
e $\card{B} \leq \card{A}$ por \olref{lem:CardinalsBehaveRight}. Logo
$\card{A} = \card{B}$ e $\cardeq{A}{B}$ por tricotomia e
\olref{lem:CardinalsBehaveRight}.
\end{proof}
\noindent
Embora seja uma prova muito simples, apoia-se implicitamente na Substituição (para assegurar
\olref[ordinals][ordtype]{thmOrdinalRepresentation}) e no
Bom ordenamento (para garantir \olref{lem:CardinalsBehaveRight}). Em
contraste, a prova na \olref[sfr][infinite][card-sb]{sec} era muito
mais autônoma (de fato, pode efetuar-se em~$\Zminus$).

\end{document}
